\documentclass[aspectratio=169]{beamer}

\usetheme{Madrid}
\usecolortheme{seahorse}
\usepackage{booktabs}
\usepackage{graphicx}
\usepackage{array}
\usepackage{hyperref}

\title[LoCoMo Compression Experiment]{LoCoMo Compression Experiment: Exp 002}
\subtitle{Budgeted retrieval, neighbor windows, semantic judging, and failure modes}
\author{Codex analysis}
\date{June 8, 2026}

\newcommand{\method}[1]{\texttt{#1}}
\newcommand{\metric}[1]{\textbf{#1}}

\begin{document}

\begin{frame}
  \titlepage
\end{frame}

\begin{frame}{What Was Checked}
\begin{itemize}
  \item Pulled latest repository state and inspected result artifacts.
  \item Found output files in \method{results/exp\_002\_merged/}:
  \begin{itemize}
    \item \method{results.parquet}
    \item \method{ablation\_table.csv}
    \item \method{category\_table.csv}
    \item \method{failure\_cases.csv}
    \item \method{plots/budget\_vs\_performance.png}
    \item \method{plots/attention\_heatmap.png}
  \end{itemize}
  \item No additional result directories, token-NLL plots, or true early/mid/late attention-map images were present.
\end{itemize}
\end{frame}

\begin{frame}{Experiment Setting}
\begin{itemize}
  \item Dataset/task: LoCoMo long-conversation QA.
  \item Sample size: 32 stratified QAs.
  \item Methods: 11 compression/retrieval methods.
  \item Budgets: 8 budget labels: 64, 128, 256, 512, 1024, 2048, 4096, and unbounded.
  \item Total rows: \(32 \times 11 \times 8 = 2816\).
  \item Model: \method{Qwen/Qwen3-4B-Instruct-2507}.
  \item Context format: quoted evidence blocks with session, date, turn id, speaker, and text.
\end{itemize}
\end{frame}

\begin{frame}{What Does ``Budget'' Mean?}
\metric{Definition.} The budget is the maximum number of tokenizer tokens allowed in the retained conversation context before asking the QA question.

\vspace{0.5em}
\begin{itemize}
  \item A small budget, e.g. 64, forces a compressor to keep only a tiny amount of memory.
  \item A larger budget, e.g. 4096, lets it keep more retrieved or recent turns.
  \item Unbounded / \method{None} means the compressor's natural output is used with no token cap.
  \item Some controls ignore budget: \method{no\_compression}, \method{oracle\_evidence}, and \method{last\_k\_turns}.
\end{itemize}

\vspace{0.5em}
\metric{Example.} For \method{dense\_retrieval}, judge score rose from 0.063 at budget 64 to 0.625 at budget 4096, showing that the cap was actually binding.
\end{frame}

\begin{frame}[fragile]{What Does ``Context Format'' Mean?}
\metric{Definition.} Context format is how retained conversation turns are written into the model prompt.

\vspace{0.5em}
\metric{Old chat-style risk.} Retrieved non-contiguous turns can look like a real dialogue history, even though turns may be missing.

\vspace{0.5em}
\metric{Exp 002 evidence format.} Retained turns are quoted as evidence blocks:
\begin{verbatim}
Conversation history:

[session=D2 date=2023-05-07 turn=D2:8 speaker=Caroline]
I researched adoption agencies today.

Question:
What did Caroline research?
\end{verbatim}
\vspace{-0.5em}
This preserves metadata such as session, date, turn id, and speaker.
\end{frame}

\begin{frame}{Concrete QA Examples Used}
\begin{tabular}{p{0.31\linewidth}p{0.20\linewidth}p{0.38\linewidth}}
\toprule
Question & Reference & What it tests \\
\midrule
What did Caroline research? & Adoption agencies & Old fact retrieval; verbose answers can be semantically correct but low token F1. \\
Who did Maria have dinner with on May 3, 2023? & her mother & Date-conditioned memory; exact answer is a person, not the date in the question. \\
Which recreational activity was James pursuing on March 16, 2022? & bowling & Context may contain the answer, but the model can still choose a distracting activity. \\
\bottomrule
\end{tabular}
\end{frame}

\begin{frame}{LoCoMo QA Categories}
\metric{Definition.} LoCoMo labels each QA with a category number describing the reasoning type. Exp 002 evaluates categories 1--4 and filters out category 5.

\vspace{0.5em}
\begin{tabular}{p{0.13\linewidth}p{0.28\linewidth}p{0.48\linewidth}}
\toprule
Category & Name & Meaning \\
\midrule
1 & Single-hop & Retrieve one fact from the conversation. \\
2 & Multi-hop & Combine multiple facts or turns. \\
3 & Temporal & Reason about dates, order, or when events happened. \\
4 & Commonsense / world knowledge & Combine conversation facts with outside/common knowledge. \\
5 & Adversarial & Information is absent; model should avoid hallucinating. Not used here. \\
\bottomrule
\end{tabular}
\end{frame}

\begin{frame}{Compression Methods}
\begin{tabular}{p{0.29\linewidth}p{0.62\linewidth}}
\toprule
Method & Logic \\
\midrule
\method{no\_compression} & Keep all turns. Budget ignored. \\
\method{last\_k\_turns} & Keep last 20 turns. Budget ignored. \\
\method{sliding\_window} & Keep most recent turns that fit budget. \\
\method{retrieval} & TF-IDF ranked turns, budget filled from candidates. \\
\method{bm25} & BM25 lexical ranked turns. \\
\method{dense\_retrieval} & Sentence embedding retrieval using MiniLM. \\
\method{*\_window} & Expand retrieved turns with neighboring turns. \\
\method{hybrid} & Recent turns plus TF-IDF retrieval. \\
\method{oracle\_evidence} & Gold evidence only; reported as upper bound. \\
\bottomrule
\end{tabular}
\end{frame}

\begin{frame}{What Does \method{*\_window} Mean?}
\metric{Definition.} A \method{*\_window} method first retrieves seed turns, then adds neighboring turns before and after each seed.

\vspace{0.5em}
\begin{itemize}
  \item \method{retrieval}: keep only TF-IDF selected turns.
  \item \method{retrieval\_window}: keep TF-IDF selected turns plus nearby turns.
  \item \method{bm25\_window}: BM25 seeds plus nearby turns.
  \item \method{dense\_retrieval\_window}: embedding-retrieved seeds plus nearby turns.
\end{itemize}

\vspace{0.5em}
\metric{Why it helps.} A gold fact is often spread across adjacent dialogue turns. The retrieved turn may mention ``adoption'', while the previous or next turn clarifies who, when, or why.
\end{frame}

\begin{frame}{What Is Dense Window?}
\metric{Dense retrieval.} Encode each conversation turn as a semantic embedding using \method{sentence-transformers/all-MiniLM-L6-v2}; encode the question the same way; rank turns by embedding similarity.

\vspace{0.5em}
\metric{Dense window.} After dense retrieval selects the most semantically similar seed turns, add neighboring turns around each seed.

\vspace{0.5em}
\begin{itemize}
  \item Base dense retrieval answers: ``Which turns are semantically closest to the question?''
  \item Dense window answers: ``Which nearby mini-dialogue around those turns should the model read?''
  \item In Exp 002, unbounded dense retrieval scored 0.609, while dense window scored 0.694.
\end{itemize}
\end{frame}

\begin{frame}{Window Example}
\begin{tabular}{p{0.18\linewidth}p{0.72\linewidth}}
\toprule
Turn & Evidence text \\
\midrule
\(D2{:}7\) & Caroline says she is thinking seriously about becoming a parent. \\
\(D2{:}8\) & Caroline says she researched adoption agencies. \\
\(D2{:}9\) & The conversation continues with paperwork and next steps. \\
\bottomrule
\end{tabular}

\vspace{0.7em}
If retrieval selects only \(D2{:}8\), it may answer the direct fact. If the question requires context, \method{retrieval\_window} keeps \(D2{:}7\), \(D2{:}8\), and \(D2{:}9\), making the retained memory more coherent.
\end{frame}

\begin{frame}{What Is the Oracle Evidence Upper Bound?}
\metric{Definition.} \method{oracle\_evidence} keeps the dataset's gold evidence turns directly. It uses labels that a real retrieval system would not know at test time.

\vspace{0.5em}
\begin{itemize}
  \item It is not a deployable compression method.
  \item It estimates how well the model could answer if retrieval were nearly perfect.
  \item In plots/tables it is labeled \method{upper\_bound\_oracle}.
  \item The gap between best real method and oracle is the remaining headroom for better retrieval, reranking, or context construction.
\end{itemize}

\vspace{0.5em}
\metric{Exp 002 example.} Oracle judge score is 0.766; best real method \method{retrieval\_window} is 0.726. The small gap suggests retrieval is close on this sample, but not solved.
\end{frame}

\begin{frame}{Experiment Logic}
\begin{enumerate}
  \item Build a stratified QA worklist from conversations.
  \item For each QA, compress the full conversation using each method and budget.
  \item Format retained context as evidence, not as fake chat history.
  \item Ask the model to answer from retained context only.
  \item Score answer quality, evidence retention, cost, latency, and perplexity diagnostics.
  \item Aggregate rows into ablation and category tables with bootstrap confidence intervals.
\end{enumerate}
\end{frame}

\begin{frame}{Metric: LLM Judge Score}
\metric{Definition.} A model-based semantic score in \(\{0, 0.5, 1\}\) judging whether the prediction answers the question equivalently to the reference.

\vspace{0.5em}
\metric{Example from this experiment.}
\begin{itemize}
  \item Question: ``What did Caroline research?''
  \item Reference: ``Adoption agencies''
  \item Prediction: ``Caroline researched adoption agencies and career options...''
  \item Token F1 is low because the prediction is verbose, but LLM judge score is 1.0 because the key answer is present.
\end{itemize}
\end{frame}

\begin{frame}{Metric: Token F1}
\metric{Definition.} Lowercase split tokens, count overlap with the reference, then compute
\[
F1 = \frac{2PR}{P+R}.
\]

\vspace{0.5em}
\metric{Example.}
\begin{itemize}
  \item Reference: ``Adoption agencies''
  \item Prediction includes many extra words: ``Caroline researched adoption agencies and career options...''
  \item Overlap exists, but precision is low because many predicted tokens are extra.
  \item Result in failure table: token F1 about 0.11 while LLM judge is 1.0.
\end{itemize}
\end{frame}

\begin{frame}{Metric: ROUGE-L}
\metric{Definition.} Longest-common-subsequence F-measure between prediction and reference.

\vspace{0.5em}
\metric{How it behaves here.}
\begin{itemize}
  \item Rewards phrase order overlap.
  \item Like token F1, it penalizes verbose correct answers.
  \item Example: ``Adoption agencies'' appears inside a long generated answer, so ROUGE-L is nonzero but still small.
\end{itemize}
\end{frame}

\begin{frame}{Metric: Exact Match}
\metric{Definition.} Normalize text by lowercasing, removing articles and punctuation, and collapsing whitespace. Score 1 only if normalized prediction equals normalized reference.

\vspace{0.5em}
\metric{Example.}
\begin{itemize}
  \item Reference: ``her mother''
  \item Prediction: ``The conversation history does not mention who Maria had dinner with on May 3, 2023.''
  \item Exact match: 0.
  \item In this run, exact match is mostly 0 because model answers are often verbose or explanatory.
\end{itemize}
\end{frame}

\begin{frame}{Metric: Date, Number, and Entity F1}
\metric{Definition.} Extract dates, numbers, and capitalized entity-like spans with regexes; compute set-level F1 against extracted reference items.

\vspace{0.5em}
\metric{Example.}
\begin{itemize}
  \item Question: ``Who did Maria have dinner with on May 3, 2023?''
  \item Reference: ``her mother''
  \item Prediction mentions ``May 3, 2023'' but not ``her mother''.
  \item Date F1 can be 0 or misleadingly about question text, so this metric is useful only when the reference itself contains a date.
\end{itemize}
\end{frame}

\begin{frame}{Metric: Evidence Recall and Precision}
\metric{Evidence recall.} Fraction of gold evidence turn ids retained:
\[
\frac{|\mathrm{kept} \cap \mathrm{gold}|}{|\mathrm{gold}|}.
\]

\metric{Evidence precision.} Fraction of retained turn ids that are gold evidence:
\[
\frac{|\mathrm{kept} \cap \mathrm{gold}|}{|\mathrm{kept}|}.
\]

\vspace{0.5em}
\metric{Example.}
If gold evidence is \(\{D2{:}8\}\) and kept turns include \(D2{:}8\), recall is 1.0. If 20 turns were kept and only one is gold, precision is 0.05.
\end{frame}

\begin{frame}{Metric: Evidence Session Recall}
\metric{Definition.} Convert gold evidence turn ids to their sessions. Score the fraction of gold evidence sessions represented in retained turns.

\[
\mathrm{EvidenceSessionRecall}
= \frac{|\mathrm{kept\ sessions} \cap \mathrm{gold\ evidence\ sessions}|}
{|\mathrm{gold\ evidence\ sessions}|}.
\]

\vspace{0.25em}
\metric{Terminology.} In this report, ``session recall'' is shorthand for \metric{evidence session recall}; it is not a separate metric.

\vspace{0.5em}
\metric{Why it matters.}
\begin{itemize}
  \item Retrieval can miss the exact gold turn but keep the right session.
  \item Example from results: \method{retrieval\_window} unbounded has evidence recall 0.585 but session recall 0.988.
  \item This means it often finds the right conversation episode even if exact turn labels are missed.
\end{itemize}
\end{frame}

\begin{frame}{Metric: Answerable Context Rate}
\metric{Definition.} Normalize the reference answer and retained text, then check whether the reference string literally appears in retained context.

\vspace{0.5em}
\metric{Example.}
\begin{itemize}
  \item Reference: ``bowling''
  \item Context contains ``bowling'', so answerable context rate is 1.0.
  \item Model prediction: ``video gaming''.
  \item Judge score is 0, showing that answerable context does not guarantee correct answer generation.
\end{itemize}
\end{frame}

\begin{frame}{Metric: Perplexity and Token NLL}
\metric{Definition.} Negative log likelihood is computed over reference answer tokens after left-truncating context so answer tokens remain scored. Perplexity is \(\exp(\mathrm{NLL})\).

\vspace{0.5em}
\metric{Use in this experiment.}
\begin{itemize}
  \item Intended as a diagnostic of how surprising the reference answer is under the retained context.
  \item Values are very large and unstable in the CSV, so they should not be used as the primary ranking metric.
  \item Token-NLL output exists in row data, but token-NLL plots were not present in pulled artifacts.
\end{itemize}
\end{frame}

\begin{frame}{Full/Unbounded Ranking}
\begin{tabular}{lrrrr}
\toprule
Method & Judge & Token F1 & Evidence Recall & Latency \\
\midrule
Upper-bound oracle & 0.766 & 0.174 & 0.990 & 0.42s \\
\method{retrieval\_window} & 0.726 & 0.114 & 0.585 & 1.91s \\
\method{dense\_retrieval\_window} & 0.694 & 0.110 & 0.763 & 2.07s \\
\method{sliding\_window} & 0.661 & 0.106 & 0.990 & 10.54s \\
\method{no\_compression} & 0.661 & 0.106 & 0.990 & 10.62s \\
\method{bm25\_window} & 0.609 & 0.092 & 0.583 & 2.17s \\
\method{last\_k\_turns} & 0.141 & 0.024 & 0.000 & 0.61s \\
\bottomrule
\end{tabular}
\end{frame}

\begin{frame}{Budget Curve: Dense Retrieval}
\begin{tabular}{rrrr}
\toprule
Budget & Judge & Evidence Recall & Latency \\
\midrule
64 & 0.063 & 0.031 & 0.41s \\
128 & 0.226 & 0.203 & 0.45s \\
256 & 0.419 & 0.266 & 0.84s \\
512 & 0.468 & 0.336 & 0.47s \\
1024 & 0.581 & 0.378 & 0.62s \\
2048 & 0.581 & 0.479 & 0.72s \\
4096 & 0.625 & 0.544 & 0.94s \\
Unbounded & 0.609 & 0.555 & 1.00s \\
\bottomrule
\end{tabular}

\vspace{0.5em}
The budget curve now slopes, unlike exp 001.
\end{frame}

\begin{frame}{Neighbor Windows Help}
\begin{tabular}{lrrr}
\toprule
Comparison & Judge Delta & Evidence Delta & Latency Delta \\
\midrule
\method{retrieval} to \method{retrieval\_window} & +0.242 & +0.250 & +1.00s \\
\method{bm25} to \method{bm25\_window} & +0.109 & +0.158 & +1.10s \\
\method{dense} to \method{dense\_window} & +0.084 & +0.209 & +1.07s \\
\bottomrule
\end{tabular}

\vspace{0.5em}
The strongest methodological result is not simply ``dense is best''; it is that retrieved turns need nearby context.
\end{frame}

\begin{frame}{Category-Level Full/Unbounded Results}
\begin{tabular}{lrrrr}
\toprule
Category & No comp. & Retrieval win. & Dense win. & Oracle \\
\midrule
1 & 0.778 & 0.722 & 0.722 & 0.944 \\
2 & 0.143 & 0.571 & 0.429 & 0.500 \\
3 & 0.643 & 0.571 & 0.571 & 0.857 \\
4 & 1.000 & 1.000 & 1.000 & 0.750 \\
\bottomrule
\end{tabular}

\vspace{0.5em}
Category variation is large, so conclusions should be stratified by QA type in the next run.
\end{frame}

\begin{frame}{Main Findings}
\begin{itemize}
  \item Budget constraints now matter: retrieval quality rises as token budget increases.
  \item \method{retrieval\_window} beats full context on judge score while being about 5x faster.
  \item Windowed retrieval is consistently better than base retrieval.
  \item Full context is not automatically best; irrelevant history appears to distract the model.
  \item \method{last\_k\_turns} performs poorly, confirming the task targets older memories.
  \item Oracle remains above real methods, so retrieval/context construction can still improve.
\end{itemize}
\end{frame}

\begin{frame}{Plot: Budget vs Token F1}
\centering
\includegraphics[height=0.76\textheight]{results/exp_002_merged/plots/split/token_f1.png}
\end{frame}

\begin{frame}{Plot: Budget vs Exact Match}
\centering
\includegraphics[height=0.76\textheight]{results/exp_002_merged/plots/split/exact_match.png}
\end{frame}

\begin{frame}{Plot: Budget vs LLM Judge Score}
\centering
\includegraphics[height=0.76\textheight]{results/exp_002_merged/plots/split/llm_judge_score.png}
\end{frame}

\begin{frame}{Plot: Budget vs Evidence Recall}
\centering
\includegraphics[height=0.76\textheight]{results/exp_002_merged/plots/split/evidence_recall.png}
\end{frame}

\begin{frame}{Plot: Budget vs Evidence Session Recall}
\centering
\includegraphics[height=0.76\textheight]{results/exp_002_merged/plots/split/evidence_session_recall.png}
\end{frame}

\begin{frame}{Plot: Budget vs Answerable Context Rate}
\centering
\includegraphics[height=0.76\textheight]{results/exp_002_merged/plots/split/answerable_context_rate.png}
\end{frame}

\begin{frame}{Plot: Evidence Recall Heatmap}
\metric{Definition.} This heatmap summarizes \metric{evidence recall} for each method and budget. Each cell is the mean fraction of gold evidence turns retained across the 32 QAs.

\vspace{0.4em}
\begin{itemize}
  \item Rows are methods.
  \item Columns are token budgets.
  \item Brighter cells mean more gold evidence turns survived compression.
  \item It is not a transformer attention map; the file name is misleading.
\end{itemize}
\end{frame}

\begin{frame}{Evidence Recall Heatmap}
\centering
\includegraphics[width=0.72\linewidth]{results/exp_002_merged/plots/attention_heatmap.png}

\vspace{0.5em}
\small Example: \method{last\_k\_turns} is dark because it keeps recent turns while the questions target older evidence; \method{upper\_bound\_oracle} is bright because it receives gold evidence directly.
\end{frame}

\begin{frame}{Statistical Caution}
\begin{itemize}
  \item Only 32 QAs were evaluated.
  \item Bootstrap confidence intervals overlap:
  \begin{itemize}
    \item \method{retrieval\_window}: 0.726, CI [0.565, 0.855]
    \item \method{no\_compression}: 0.661, CI [0.500, 0.806]
    \item oracle: 0.766, CI [0.609, 0.906]
  \end{itemize}
  \item The result is promising, not definitive.
  \item Next experiment should use at least 200 stratified QAs.
\end{itemize}
\end{frame}

\begin{frame}{Code/Reporting Issues Found}
\begin{itemize}
  \item Result artifacts include no true early/mid/late attention example PNGs.
  \item \method{attention\_heatmap.png} is misnamed; it is not a transformer attention map.
  \item Token-NLL plots were expected by code but are absent from pulled outputs.
  \item Non-budgeted methods are plotted across all budgets, which can mislead budget comparisons.
  \item \method{failure\_cases.csv} is very large because token F1 flags many semantically correct verbose answers as failures.
\end{itemize}
\end{frame}

\begin{frame}{Metric Issues and Improvements}
\begin{itemize}
  \item LLM judge uses the same model family as answer generation; use an external stronger judge for final claims.
  \item Literal \metric{answerable context rate} misses paraphrases and aliases.
  \item Regex entity/date extraction is brittle; use a date parser and NER if this becomes a core metric.
  \item Perplexity values are huge and unstable; keep token NLL as diagnostic, not headline ranking.
  \item Exact match is too strict for verbose model answers; use concise decoding prompts or postprocess answers.
\end{itemize}
\end{frame}

\begin{frame}{Method Improvements}
\begin{itemize}
  \item Increase sample size and report paired bootstrap deltas.
  \item Tune neighbor window size by category and budget.
  \item Retrieve at session-chunk level, then compress within the retrieved session.
  \item Add reranking after dense/BM25 retrieval.
  \item Separate budget-sensitive methods from fixed controls in all plots.
  \item Add true attention example export or rename the existing heatmap to avoid confusion.
\end{itemize}
\end{frame}

\begin{frame}{Bottom Line}
\begin{block}{Best current interpretation}
Small, relevant, neighbor-expanded memory is a better latency/accuracy tradeoff than dumping the full LoCoMo conversation into the model.
\end{block}

\begin{block}{Confidence level}
The conclusion is directionally strong but statistically preliminary because the run uses only 32 QAs and the judge-score confidence intervals overlap.
\end{block}
\end{frame}

\end{document}
